\documentclass[main.tex]{subfiles}


\begin{document}

% \AddToShipoutPictureFG{%
% \begin{tikzpicture}[overlay,remember picture]
% \path (current page.south west) -- (current page.north east)
% node[midway,scale=1,color=gray,sloped,opacity=0.4] {\textnormal{Кравченко Игорь Игоревич\hspace{1cm} +79010144910\hspace{1cm} super.antig@yandex.ru}};
% \end{tikzpicture}
% }


% %from pagenumber to up
% \phantomsection 
% \hypertarget{toc}{}

 

 

% номер секции вручную
\setcounter{section}{19
}
\addtocounter{section}{-1} 

\section{Спутник}

\textbf{Спутник}~--- это тело, двигающееся вокруг другого тела под действием силы тяготения. 

Можно сказать, что спутником, например, Солнца является любая планета Солнечной системы, в том числе Земля. Земля, в свою очередь, также имеет собственный \emph{естественный спутник}~--- Луну. 

% В этих примерах речь идет о так называемых \emph{естественных спутниках}.

На рис.\,\ref{fig:p37} условно показано круговое движение спутника вокруг планеты.

    \begin{figure}[H]
    \centering

    \begin{tikzpicture}[font=\small,            line cap=round,                             >={Stealth[inset=0pt,round,length=3mm, angle'=20]},vec/.style={>={Stealth[inset=0pt,round,length=3.5mm, angle'=20]}}]


\shade[ball color=Green] (0,0) circle (.6);
\node [below,black] at (0,-.6) {П};


\draw[dashed] (0,0) circle (3cm);
\path (0,0) ++ (20:3cm) node (s) {} ++ (.15,0) node[black,right] {С};
\shade[ball color=lightgray] (s) circle (.15);


%vec
\draw[->,NavyBlue,thick,vec] (s.center) --++ (110:2cm) node[black,right,midway] {$\vec v_\text{к}$};
\node at (s) [circle,fill=red,inner sep=0pt,minimum size=1mm,label=below:] {};
\draw[->,red,thick,vec] (s.center) --++ (200:1cm) node[black,below] {$\vec F_\text{т}$};

%r
\draw[<->,gray] (0,0) --++ (90:3cm) node[midway,left,black] {$r$};
\node at (0,0) [circle,fill=gray,inner sep=0pt,minimum size=1mm,label=below:] {};


    \end{tikzpicture}
\caption{Вращение спутника вокруг планеты}
\label{fig:p37}
\end{figure}

Спутник~С, обладающий касательной скоростью~$\vec v_\text{к}$, под действием силы тяготения~$\vec F_\text{т}$, направленной к центру планеты~П, вращается вокруг планеты, как говорят, по круговой орбите~--- то есть по окружности (штриховая линия)\footnote{В действительности траекториями движения спутников (орбитами) являются \emph{эллипсы}.}. Центр этой окружности совпадает с центром планеты, расстояние~$r$ между центрами тел (радиус орбиты) не меняется. Центростремительное ускорение спутника по сути является ускорением свободного падения в месте нахождения спутника и создается силой тяготения планеты ($a_\text{ц}=g$).

\textbf{Искусственный спутник}~--- это спутник, созданный человеком. Для того, чтобы такой спутник смог двигаться по круговой орбите вокруг планеты \emph{вблизи ее поверхности} (то есть высота над поверхностью $h=0$), он должен обладать так называемой \textbf{первой космической скоростью}
\begin{equation}
    v_\text{1к}=\sqrt{G\dfrac{m_\text{планеты}}{R_\text{планеты}}},
\end{equation}
где~$m_\text{планеты}$ и~$R_\text{планеты}$~--- масса и радиус планеты соответственно.

% \textbf{Искусственный спутник}~--- это спутник, созданный человеком. 

Если скорость запущенного спутника меньше первой космической скорости, то он упадет на планету. При скорости, большей первой космической, спутник опишет эллиптическую траекторию. 

Минимальная скорость, требуемая запущенному телу для удаления от планеты на бесконечность, называется \textbf{второй космической скоростью} и вычисляется по формуле:
\begin{equation}
    v_\text{2к}=\sqrt{2}v_\text{1к}.
\end{equation}

Так, для Земли первая космическая скорость приблизительно равна~$8$~км$/$с.

% $v_\text{1к}\approx8$~км$/$с.

 
\end{document}